\documentclass[11pt,a4paper]{article}
\usepackage[T1]{fontenc}
\usepackage{lmodern,amsmath,amssymb,booktabs,microtype}
\usepackage[margin=24mm]{geometry}
\usepackage[colorlinks=true,urlcolor=blue,linkcolor=blue]{hyperref}
\newcommand{\Js}{\mathcal J_s}
\newcommand{\ket}[1]{|#1\rangle}
\newcommand{\norm}[1]{\lVert#1\rVert}
\title{Two-singlet channels and the first interaction signal\\A chord-space resolution of the short-time commutator}
\author{Alan Day \and Chengheng\thanks{AI research assistant (OpenAI).}}
\date{Local research note --- 6 November 2026}
\begin{document}\maketitle
\begin{abstract}
We resolve the previously computed singlet commutator curvature into
orthogonal chord channels. Four fifths of its leading small-coupling
coefficient comes from the first relative two-matter channel, whose
same-species norm vanishes in the free limit; one fifth comes from a
zero-matter component proportional to the Hamiltonian. The exact
finite-parameter sum agrees with an independent full Wick enumeration.
We also extend the positive chord generator to ordered two-matter
representations and give their standard Jacobi--Laguerre limiting
realization. This separates relative-channel labels from physical masses
and identifies the extra operator data needed for a meson spectrum.
\end{abstract}

\section{An exact decomposition of an existing observable}
The state is the full annealed normalized trace $\tau$, with
$\beta=\mu_{\rm chem}=0$. Take complex DSSYK double scaling first,
then small $h=4\lambda$, $q=e^{-h}$, at fixed times or fixed chord
degrees. Our normalized conserved-subtracted derivative singlet $X$
has $d=2$, $\tau X^2=1$, mixed crossing weight $r=q^2$, and
$X$--$X$ crossing weight $\chi=q^4=r^2$. Write
\begin{equation}
 H_s=2\Js H_{\rm src},\qquad \tau H_{\rm src}^2=h^{-1},\qquad
 F=H_{\rm src}X^2-2XH_{\rm src}X+X^2H_{\rm src}.
\end{equation}
Since $[\dot X,X]=2i\Js F$, the nonnegative commutator square is
\begin{equation}
 \mathcal C(t)=\tau([X(t),X]^\dagger[X(t),X])
 =4\Js^2t^2\norm F^2+O(t^4).
\end{equation}
Our channel calculation gives the exact fixed-$h$ decomposition
\begin{equation}
 \boxed{\norm F^2=\frac1h\left[
 \underbrace{4(1-r)^2}_{\text{zero matter}}
 +\underbrace{2(1-r)^4}_{k=0\ \text{two matter}}
 +\underbrace{4(1-r^2)^2}_{k=1\ \text{two matter}}\right].}\label{result}
\end{equation}
Thus the respective leading coefficients in $\mathcal C(t)$ are
$256\Js^2\lambda t^2$, zero at order $\lambda$, and
$1024\Js^2\lambda t^2$. Their sum reproduces the earlier
$1280\Js^2\lambda t^2$. This resolves the existing scalar result
into states; it is not a new Lyapunov exponent.

\newpage
\section{Proof using the two-open-matter Gram form}
Let $(L,M,R)$ denote one Hamiltonian chord respectively before, between,
or after two open $X$ chords. The no-crossing and crossed matter
pairings give
\begin{equation}
 G=\begin{pmatrix}1&r&r^2\\r&1&r\\r^2&r&1\end{pmatrix},\qquad
 G_\times=\chi\begin{pmatrix}1&r&r^2\\r&r^2&r\\r^2&r&1\end{pmatrix}.
 \label{gram}
\end{equation}
The middle crossed entry is $\chi r^2$, not $\chi$; the Hamiltonian
chord crosses both matter chords. The physical two-open-$X$ norm is
$G+G_\times$. Internal matter contractions form the separate
zero-matter sector. This is the same-species Wick-sector decomposition,
not the assertion that $G_\times$ is a unitary exchange matrix.

In the ordered direct form $G$, the normalized first primary and the
symmetric external descendant are
\begin{equation}
 u_1=\frac{(-r,1+r^2,-r)^T}{\sqrt{1-r^4}},\qquad
 s_0=(1,0,1)^T.
\end{equation}
Both external annihilators kill $u_1$; it is the $d_{\rm eff}=5$
primary. Direct multiplication proves
\begin{gather}
 u_1^TGu_1=1,\quad u_1^TG s_0=u_1^TG_\times s_0=0,\nonumber\\
 u_1^TG_\times u_1=-\chi r^2,\quad
 s_0^TG s_0=2(1+r^2),\quad s_0^TG_\times s_0=2\chi(1+r^2).
\end{gather}
The raw two-matter part of $\sqrt hF$ has vector $v=(1,-2,1)^T$.
Its exact projection is
\begin{equation}
 v=c\,s_0+a\,u_1,\qquad
 c=\frac{(1-r)^2}{1+r^2},\qquad
 a=-\frac{2(1-r^2)}{\sqrt{1-r^4}}.\label{project}
\end{equation}
Substituting $\chi=r^2$ yields
\begin{align}
 \norm{c s_0}_{G+G_\times}^2&=2(1-r)^4,\nonumber\\
 \norm{a u_1}_{G+G_\times}^2&=a^2(1-r^4)=4(1-r^2)^2.
\end{align}
The contraction of the two $X$'s inside $H_{\rm src}X^2$ or
$X^2H_{\rm src}$ leaves $H_{\rm src}$; inside $XH_{\rm src}X$
it leaves $rH_{\rm src}$. The zero-matter part is therefore
$2(1-r)H_{\rm src}$, with squared norm $4(1-r)^2/h$.
Orthogonality of the zero- and two-open-matter Wick sectors proves
Eq.~\eqref{result}.

The $k=1$ direct unit vector has physical squared norm
$1-r^4=1-e^{-8h}\sim8h$. Simultaneously $a^2\sim8h$.
Their product divided by $h$ gives an order-$h$ effect in $\norm F^2$.
Dropping this channel because its unrescaled free-limit norm vanishes
would discard the leading $4/5$ contribution.

\newpage
\section{Ordered two-matter representations and the positive generator}
For dimensions $d_1,d_2>0$, the ordered chord states are
$\ket{l,m,r}$, with $E=\bar n=d_1+d_2+l+m+r$. The coproduct gives
\begin{align}
 a_L\ket{l,m,r}={}&[l]_q\ket{l-1,m,r}
 +q^{d_1+l}[m]_q\ket{l,m-1,r}\nonumber\\
 &+q^{d_1+d_2+l+m}[r]_q\ket{l,m,r-1},\label{action}
\end{align}
and the reflected $a_R$ action. We use the physical Gram form, whose
pure-middle direct norm is $[m]_q!$. These are the existing
Lin--Stanford rules \cite{LS}, Sec.~3.3 and App.~B.2.

Their decomposition is
\begin{equation}
 \mathcal R_{2}^{d_1,d_2}=\bigoplus_{k=0}^\infty\mathcal R_1^{D_k},
 \qquad D_k=d_1+d_2+k.\label{irreps}
\end{equation}
One may construct its normalized primary $\Omega_k$ without an
all-$k$ normalization formula. At grade $k$, take the orthogonal
complement of the external creations from grade $k-1$. Adjointness
then gives $a_L\Omega_k=a_R\Omega_k=0$. Inductively the already
constructed families fill all but one dimension, since
\begin{equation}
 \dim\mathcal H_k=\binom{k+2}{2},\qquad
 \sum_{j=0}^{k-1}(k-j+1)=\binom{k+2}{2}-1.
\end{equation}
Positivity of the chord Gram form and the one-matter irreps ensure
nondegenerate descendant blocks and orthogonality of different
families. This gives one new primary at each grade.

Define $K_+=(a_L^\dagger-a_R^\dagger)/\sqrt h$ and
$K_-=(a_L-a_R)/\sqrt h$ as in the preceding note, and
\begin{equation}
 B_h=(K_++K_-)/2,\qquad \Pi_h=(K_--K_+)/(2i),\qquad Q_h=E+\Pi_h.
\end{equation}
Each summand in Eq.~\eqref{irreps} has the preceding one-matter
proof with $d=D_k$. Therefore $Q_h$ is positive and self-adjoint
on the ordered two-matter space (domain $D(E)$ at fixed $h$).
For every fixed $k,n$ its normalized antisymmetric descendants tend
to $D^+_{D_k}$, with raising coefficient
$\sqrt{(n+1)(2D_k+n)}$.

The resulting contracted sector is $\bigoplus_k D^+_{D_k}$ on
finite sums of fixed-channel descendants. It is not a uniform
limit of all finite-$h$ chord directions. In particular it does not
retain every singularly rescaled ordering difference automatically.
The original concatenation coproduct is not a primitive sum on
independent matter Hilbert spaces. Xu's factorization likewise
requires a modified inner product \cite{Xu}.

\newpage
\section{The relative coordinate and why $k$ is not a mass}
The standard positive discrete-series tensor decomposition
\cite{KV,J} realizes Eq.~\eqref{irreps} after contraction as
$D^+_{d_1}\otimes D^+_{d_2}$. An explicit unitary coordinate map is
\begin{equation}
 P=p_1+p_2>0,\quad x=p_1/P\in(0,1),\qquad
 (\mathcal U f)(P,x)=\sqrt P\,f(Px,P(1-x)).
\end{equation}
The measure is $dP\,dx$. Let $D=d_1+d_2$, and define
\begin{align}
 \chi_k(x)&=\frac{x^{d_1-1/2}(1-x)^{d_2-1/2}}{\sqrt{N_k}}
 P_k^{(2d_2-1,2d_1-1)}(2x-1),\nonumber\\
 N_k&=\frac{\Gamma(k+2d_1)\Gamma(k+2d_2)}
 {(2k+2D-1)k!\Gamma(k+2D-1)}.\label{jacobi}
\end{align}
At removable singular parameter values, take the continuous limit
of the displayed normalization. The $\chi_k$ form an orthonormal
basis of $L^2(0,1;dx)$. With the preceding Laguerre radial basis,
the channel map is
\begin{equation}
 e_{k,n}\longmapsto\phi_n^{D_k}(P)\chi_k(x),\qquad
 Q=P,\quad B=i(P\partial_P+\tfrac12).
\end{equation}
It intertwines $E,B,Q$ with the sums of their one-particle
realizations. In particular, if $a=d_1-1/2$, $b=d_2-1/2$, the
transformed sum $S_1+S_2$ is
\begin{gather}
 -P\partial_P^2-\partial_P+\frac{\mathcal T_x}{P},\qquad
 \mathcal T_x=-\partial_x x(1-x)\partial_x+\frac{a^2}{x}
 +\frac{b^2}{1-x}+\frac14,\nonumber\\
 \mathcal T_x\chi_k=(D_k-\tfrac12)^2\chi_k.\label{T}
\end{gather}
This follows from the Jacobi differential equation. It verifies the
tensor realization explicitly, rather than identifying a finite-$h$
coproduct with a primitive sum.

For a free positive-energy extension with equal particle masses $m$
and common momentum scale $\mu$, light-front kinematics gives
\begin{equation}
 P^+=\mu P,\qquad P^-=\frac{m^2}{\mu P x(1-x)},\qquad
 M^2=\frac{m^2}{x(1-x)}.
\end{equation}
At $d_1=d_2=2$, direct integration in the \emph{relative} space gives
\begin{equation}
 \langle\chi_0|M^2/m^2|\chi_0\rangle=\frac{14}{3},\qquad
 \boxed{\langle\chi_0|M^2/m^2|\chi_2\rangle=\frac{\sqrt{11}}3\ne0.}
\end{equation}
Thus even the free mass operator mixes the discrete $k$ channels.
Its spectrum is the continuum $[4m^2,\infty)$, not the list $D_k$.
The $p_i,x$ here parametrize the two-singlet representation; they
have not been identified as quark momenta inside one meson.

\newpage
\section{Higher channels, verification and the dynamical question}
The opposite-order overlap suggested in Lin--Stanford App.~F is
\begin{equation}
 R_k=(-1)^k q^{d_1d_2+k(d_1+d_2)+k(k-1)/2}.
 \label{R}
\end{equation}
The authors report checks through $k=25$, rather than an all-$k$
proof of this formula. We use exact direct calculations for $k=0,1$
in Eq.~\eqref{result}; our further tests through $k=4$ do not upgrade
Eq.~\eqref{R} to a theorem. Its exponent can be written
\begin{equation}
 \tfrac12\{D_k(D_k-1)-d_1(d_1-1)-d_2(d_2-1)\}.
\end{equation}
This is an overlap between differently ordered states. Its modulus
is below one at finite $h$ and it must not be interpreted as an
elastic scattering phase. For identical slots at the verified first
grade, the physical form $G+G_\times$ retains the norm defect that
produces the commutator contribution on page 2.

\paragraph{Independent checks.}
The saved validation contains 38 exact identities: first-grade
primary/projection formulas, the complete six-letter Wick sum for
$F^2$, Jacobi orthogonality, and relative Casimir eigenfunctions.
Three additional symbolic differential identities verify the full
two-variable $S$ and $B$ transformations and the Jacobi conjugation
for general parameters. A separate recursive Gram implementation
constructs primaries through $k=4$ at
\begin{equation}
 (h,d_1,d_2)=(0.05,2,2),\ (0.01,2,2),\ (0.02,1.3,2.1).
\end{equation}
Computations at 85 and 115 decimal digits differ relatively by at
most $2.5\times10^{-78}$ in the recorded overlaps and projections.
The maximum 115-digit overlap residual is below $1.8\times10^{-108}$.
These are finite-degree numerical checks, not certified error bounds
or evidence for uniform large-$k$ limits.

\paragraph{The remaining physical input.}
In the tensor realization $Q$ and $B$ act only on $P$. Any suitable
positive self-adjoint relative operator $\mathcal M_x^2$ therefore
defines kinematically
\begin{equation}
 P^+=\mu P,\qquad P^-=\frac{\mathcal M_x^2}{\mu P}.
\end{equation}
The algebra alone permits either continuous or discrete relative
spectra. To choose $\mathcal M_x^2$ requires a microscopic spatial
dictionary or independent dynamical Ward identities; channel counting
does not choose it. The computed commutator decomposition supplies an
interaction constraint that a candidate must reproduce.

There is also a distinction of sectors: two open $X$ chords describe
products of our \emph{single} derivative singlet. They are useful for
singlet interactions. The single-meson tower proposed by
Miyashita--Sekino--Susskind instead calls for independent derivative
fermion bilinears \cite{MSS}. Their spectral matrix is the next
direct microscopic target; the present relative index must not be
relabeled as that tower.

\newpage
\section{Next calculation and reproducibility}
We now have a positive total chord generator and an exact first
relative-channel contribution to the noncommutativity of $X$.
The next task is to enlarge the microscopic singlet operator family
beyond $X$, so that particle-species information can be distinguished
from time derivatives and products of the same singlet.

Start with the next derivative bilinears
$\sum_i c_i^\dagger\operatorname{ad}_{H_s}^{n}(c_i)$ and their
adjoints. Resolve their Hermitian components, conserved projections,
and Liouvillian descendants before computing the remaining independent
Gram and spectral matrix. The same disorder couplings must be retained
inside the derivatives. Use existing complex-SYK chord/embedded-ensemble
identities where available, followed by exact low-order checks. The
criterion is a nonzero new independent spectral direction, not merely
another correlator of $X$.

The channel result here remains a separate constraint for the eventual
interaction map. We have not equated it with a spatial on-shell
amplitude. The ordinary fixed-$k,n$ contraction also leaves open
singular ordering sectors, growing times, and the small-$P$ domain of
inverse translations. Those questions can be addressed when a concrete
observable requires them, rather than replacing the spectral calculation.

\paragraph{Saved files.}
The stage \texttt{susskind-two-matter} contains the LaTeX source,
Chinese derivation, literature map, numerical outputs and SHA-256
manifest. The numerical code is \nolinkurl{two_matter_chords.py}.
The independent checks are \nolinkurl{check_two_matter_chords.py}
and \nolinkurl{check_two_matter_tensor.py}.
The build script uses two \texttt{pdflatex} passes, and all six
rendered pages are inspected. This stage is saved locally without
a Git commit, push, or archive.

\begin{thebibliography}{9}
\bibitem{LS} H. W. Lin and D. Stanford,
\emph{A symmetry algebra in double-scaled SYK},
\href{https://arxiv.org/abs/2307.15725}{arXiv:2307.15725v1}.
Secs.~2.3, 3.3; Apps.~B.2 and F.
\bibitem{Xu} J. Xu,
\emph{On Chord Dynamics and Complexity Growth in Double-Scaled SYK},
\href{https://arxiv.org/abs/2411.04251}{arXiv:2411.04251v3}, Sec.~3.2.
\bibitem{KV} H. T. Koelink and J. Van der Jeugt,
\emph{Convolutions for orthogonal polynomials from Lie and quantum
algebra representations},
\href{https://arxiv.org/abs/q-alg/9607010}{arXiv:q-alg/9607010},
Corollary 3.6 and Remark 3.7.
\bibitem{J} J. Van der Jeugt and R. Jagannathan,
\emph{Realizations of $su(1,1)$ and $U_q(su(1,1))$ and generating
functions for orthogonal polynomials},
\href{https://arxiv.org/abs/math-ph/9807019}{arXiv:math-ph/9807019}, Eq.~(2.10).
\bibitem{MSS} S. Miyashita, Y. Sekino and L. Susskind,
\emph{Holograms and Standard Models},
\href{https://arxiv.org/abs/2607.05678}{arXiv:2607.05678v2}, Secs.~4.2--4.3.
\end{thebibliography}
\end{document}
